%% uhf_faithful_trace_glimm_2026-07-23.tex
%%
%% Full paper (expansion of the 2026-07-20 ITP note): a machine-checked
%% construction of the base-3 UHF C*-algebra M_{3^infty} with a canonical
%% faithful tracial state, and a formalized Glimm-simplicity theorem, in
%% Lean 4 over mathlib.
%%
%% All mathematical claims below are either (a) kernel-verified Lean 4
%% theorems at HEAD 91cc1a01 (2026-07-23), axioms [propext,
%% Classical.choice, Quot.sound], zero project axioms, zero sorries, or
%% (b) standard cited literature. Verified build data: lake build PF =
%% 4472 jobs, lake build (full) = 8930 jobs, mathlib pin eed770a4,
%% toolchain leanprover/lean4:v4.24.0-rc1.

\documentclass[11pt,a4paper]{amsart}

\usepackage[utf8]{inputenc}
\usepackage[T1]{fontenc}
\usepackage{lmodern}
\usepackage{amsmath,amssymb,amsthm}
\usepackage{mathtools}
\usepackage{microtype}
\usepackage{geometry}
\geometry{margin=1in}
\usepackage[hidelinks,colorlinks=false]{hyperref}
\usepackage{enumitem}
\usepackage{booktabs}
\usepackage{framed}
\usepackage{xcolor}
\usepackage{seqsplit}

\emergencystretch=8em
\hbadness=10000
\setlength{\hfuzz}{30pt}

\theoremstyle{plain}
\newtheorem{theorem}{Theorem}[section]
\newtheorem{lemma}[theorem]{Lemma}
\newtheorem{proposition}[theorem]{Proposition}
\newtheorem{corollary}[theorem]{Corollary}

\theoremstyle{definition}
\newtheorem{definition}[theorem]{Definition}
\newtheorem{remark}[theorem]{Remark}

\sloppy
\tolerance=10000
\pretolerance=10000
\emergencystretch=15em

\newcommand{\lt}[1]{\texttt{\seqsplit{#1}}}
\newcommand{\Tinf}{T_\infty}
\newcommand{\tuhf}{\tau_{\mathsf{UHF}}}
\newcommand{\opnorm}{\|\cdot\|_{\mathsf{op}}}

\title[A Machine-Checked UHF Algebra with a Faithful Trace]{A Machine-Checked Construction of the UHF Algebra \texorpdfstring{$M_{3^\infty}$}{M(3\textasciicircum\ensuremath{\infty})} with a Faithful Tracial State: Formalizing Glimm Simplicity in Lean~4}

\author{Pablo Cohen}
\address{Mesa, Arizona, USA}
\email{psolorzano@gmail.com}
\urladdr{ORCID:~\href{https://orcid.org/0009-0002-0734-5565}{0009-0002-0734-5565}}

\date{July 23, 2026}

\begin{document}

\begin{abstract}
\noindent We report a complete Lean~4 formalization, over
mathlib~\cite{mathlib2020,moura-ullrich2021}, of the uniformly hyperfinite
(UHF) $C^*$-algebra $M_{3^\infty}$ --- realized as the metric completion
$\Tinf$ of the inductive limit of the matrix tower $M_{3^k}(\mathbb{C})$
under the unital $*$-embeddings $A \mapsto A \otimes I_3$ --- together with
a canonical tracial state $\tuhf$ on $\Tinf$ that is kernel-verified to be
additive, unital, $1$-Lipschitz, tracial, Hermitian, positive, and
\emph{faithful}: $\tuhf(x^*x) = 0 \Rightarrow x = 0$ for every $x \in \Tinf$.
As a corollary we obtain a machine-checked proof that $\Tinf$ is
algebraically simple --- every two-sided ideal is $\bot$ or $\top$. We
further prove that $\tau_{\mathsf{UHF}}$ is the \emph{unique} continuous
unital tracial functional on $\Tinf$ (theorem
\lt{substrate\_UHF\_trace\_unique}; positivity is not required). Faithfulness,
simplicity, and uniqueness of the trace together identify $\Tinf$ as the
Glimm $3^\infty$ UHF factor. To the best of our knowledge, after searching
the major proof-assistant libraries (mathlib, the Coq \texttt{opam} archive,
and the Isabelle Archive of Formal Proofs) and the formalization literature,
this is the first formalization of Glimm's 1960 simplicity theorem for a UHF
algebra~\cite{glimm1960} and the first machine-checked construction of an
infinite-dimensional simple $C^*$-algebra with a canonical faithful and
unique trace; we have not performed an exhaustive survey and would welcome
pointers to prior art.

\smallskip

\noindent The completion-tier faithfulness proof is the heart of the paper.
It routes through a trace-preserving conditional expectation
$E_k \colon \Tinf \to M_{3^k}(\mathbb{C})$ (the partial trace over the
tail), proved to be a $\|\cdot\|_2$-contraction by an elementary
Cauchy--Schwarz argument --- where the classical route would invoke the
Kadison--Schwarz inequality and complete positivity, neither of which is
available in mathlib --- and an operator-norm contraction via the isometry
decomposition $E_k B = \tfrac{1}{3}\sum_{t} W_t^* B\, W_t$. Supporting
machinery, novel as formalized content, includes a Cauchy--Schwarz
inequality for the trace pairing, the trace-null set as a mathlib-native
two-sided ideal, a discrete Weyl clock--shift unitary-averaging engine
$\tfrac{1}{n^2}\sum_{a,b}(C^aS^b)\,x\,(C^aS^b)^* = (\operatorname{tr}x/n)\cdot 1$,
algebraic simplicity of the pre-completion tower, and a four-fold
equivalence (trace-detection $\Leftrightarrow$ triviality of the trace-null
ideal $\Leftrightarrow$ faithfulness $\Leftrightarrow$ simplicity)
that precisely located the difficulty before it was overcome. Three
standalone API files of mathlib-PR quality (topological closures of
two-sided ideals; the non-unital continuous functional calculus mapping
into closed ideals; spectral projections from clopen quasispectral sets)
are contributed as by-products. Every headline theorem reports kernel axioms
$[\texttt{propext}, \texttt{Classical.choice}, \texttt{Quot.sound}]$, with
zero project axioms and zero \texttt{sorry}s, verified at pinned toolchain
and mathlib revisions.
\end{abstract}

\maketitle

\tableofcontents

%% =====================================================================
\section{Introduction}\label{sec:intro}
%% =====================================================================

Uniformly hyperfinite (UHF) $C^*$-algebras are among the oldest and most
completely understood objects of operator algebra theory. Introduced by
Glimm in 1960~\cite{glimm1960} --- with roots in the hyperfinite factor
constructions of Murray and von Neumann~\cite{murray-vonneumann1936} ---
they are the norm-closures of increasing unions of full matrix algebras,
classified completely by a supernatural number, each simple, each carrying
a unique tracial state which is automatically faithful. Standard
references establish these facts~\cite{davidson1996, blackadar2006,
takesaki1979}. None of them, before the present work, had been
machine-checked.

The gap is not accidental. Formalized operator algebra is young: mathlib
--- the Lean~4 mathematical library~\cite{mathlib2020,moura-ullrich2021}
--- provides a $C^*$-algebra typeclass hierarchy, an $L^2$-operator-norm
$C^*$-structure on matrix algebras, ring direct limits, metric completions,
and a continuous functional calculus, largely following the formalized
functional analysis program of Dupuis, Lewis, and
Macbeth~\cite{dupuis-lewis-macbeth2022}. But it contains no concrete
infinite-dimensional simple $C^*$-algebra, no tracial state theory, no
conditional expectations, no completely positive maps, and no spectral
projections. Constructing a UHF algebra and proving the classical facts
about it therefore requires building, from mathlib's primitives, a
substantial amount of infrastructure that the pen-and-paper literature
takes for granted --- and, at several points, finding proofs
\emph{different from} the classical ones, because the classical proofs
lean on theory (Kadison--Schwarz, complete positivity, Borel functional
calculus) that does not yet exist in any proof assistant library.

\subsection{Contributions}\label{sec:contributions}

This paper describes a Lean~4 formalization, kernel-verified with axioms
$[\texttt{propext}, \texttt{Classical.choice}, \texttt{Quot.sound}]$ only,
zero project-level axioms, and zero \texttt{sorry}s, of:

\begin{enumerate}[label=\textbf{(C\arabic*)}, leftmargin=3em, itemsep=3pt]
\item \textbf{The UHF algebra $M_{3^\infty}$, mathlib-native.} The
  base-3 matrix tower $A_k := M_{3^k}(\mathbb{C})$ with the unital
  $*$-embeddings $A \mapsto \mathrm{reindex}(A \otimes I_3)$, its ring
  direct limit, and the metric completion $\Tinf$ of that limit under the
  $L^2$-operator norm, carrying \texttt{Ring}, \texttt{StarRing},
  \texttt{NormedRing}, and \texttt{CStarAlgebra} instances
  (\S\ref{sec:tower}).
\item \textbf{A canonical tracial state $\tuhf$ on $\Tinf$}, obtained as
  the unique uniformly continuous extension of the normalised matrix
  traces, and kernel-verified to be additive, unital, $1$-Lipschitz,
  tracial, Hermitian, and positive (\S\ref{sec:trace}).
\item \textbf{Unconditional faithfulness of $\tuhf$ at every tier} ---
  matrix level, direct limit, and, the hard case, the completion:
  $\tuhf(x^*x) = 0 \Rightarrow x = 0$ for all $x \in \Tinf$
  (theorem \lt{UHF\_trace\_faithful}, \S\ref{sec:condexp}).
\item \textbf{Glimm simplicity of $\Tinf$}: every two-sided ideal of
  $\Tinf$ is $\bot$ or $\top$ (theorem
  \lt{substrate\_completion\_simple\_unconditional},
  \S\ref{sec:condexp}). To the best of our knowledge --- after searching
  the major proof-assistant libraries and literature (see the priority
  note in \S\ref{sec:related}) --- this is the first formalized simplicity
  theorem for an infinite-dimensional $C^*$-algebra.
\item \textbf{Uniqueness of the trace}: $\tau_{\mathsf{UHF}}$ is the unique
  continuous unital tracial functional on $\Tinf$ (theorem
  \lt{substrate\_UHF\_trace\_unique}, \S\ref{sec:condexp}); the argument
  requires only linearity, the trace identity, unitality, and continuity,
  not positivity. Together with faithfulness and simplicity this identifies
  $\Tinf$ as the Glimm $3^\infty$ UHF factor
  (\lt{r113\_substrate\_UHF\_factor\_capstone}).
\item \textbf{The proof machinery itself}, novel as formalized content: a
  Cauchy--Schwarz inequality for the trace sesquilinear pairing; the
  trace-null set as a mathlib-native \texttt{TwoSidedIdeal}; a discrete
  Weyl clock--shift unitary-averaging engine; algebraic simplicity of the
  pre-completion tower; a completion-tier averaging operator; a four-fold
  equivalence between trace-detection, ideal-triviality, faithfulness, and
  simplicity; and a trace-preserving conditional expectation
  $E_k \colon \Tinf \to M_{3^k}(\mathbb{C})$ with elementary
  $\|\cdot\|_2$- and operator-norm contraction proofs
  (\S\S\ref{sec:engine}--\ref{sec:condexp}). Three standalone
  mathlib-PR-quality API files are contributed as by-products
  (\S\ref{sec:formalization}).
\end{enumerate}

\subsection{Why faithfulness is the hard part}\label{sec:why-hard}

Everything about $\tuhf$ \emph{except} faithfulness transfers from the
matrix level to the completion by soft means. Additivity, traciality,
Hermitian symmetry, and positivity are all \emph{closed conditions}: each
asserts that a continuous function vanishes (or is nonnegative) on
$\Tinf$, so it suffices to verify it on the dense image of the direct
limit and invoke completion induction on a closed set. Matrix-level
verification is elementary, and mathlib's
\lt{UniformSpace.Completion.induction\_on} closes each one.

Faithfulness is different in kind. The statement
$\tuhf(x^*x) = 0 \Rightarrow x = 0$ is \emph{not} a closed condition ---
the set of $x$ satisfying an implication with a nontrivial conclusion is
not the preimage of a closed set --- so density arguments do not apply. A
limit of matrices, each individually detected by its trace, could in
principle converge to a nonzero element whose trace-of-square vanishes:
ruling this out requires genuinely global structure. Classically one
appeals to either (a) simplicity of the algebra --- a state whose left
kernel generates a proper closed two-sided ideal must be faithful when no
such ideal exists --- with simplicity supplied by Glimm's
theorem~\cite{glimm1960}; or (b) the norm-one conditional expectations
onto the finite levels, in the tradition of
Tomiyama~\cite{tomiyama1957}, whose contractivity is classically a
consequence of complete positivity and the Kadison--Schwarz
inequality~\cite{kadison1952}. Neither route existed in mathlib. The
formalization ultimately built route (b) from scratch, with elementary
replacements for the missing complete-positivity theory
(\S\ref{sec:condexp}), and obtained route (a)'s conclusion --- simplicity
--- as a corollary rather than an input, inverting the classical
dependency.

\subsection{An honest development history}\label{sec:history}

We describe the development arc candidly, because we believe the
trajectory is itself a datum about formalizing operator algebras, and
because an earlier version of this material contained a claim we had to
retract before it was eventually proved.

The construction of $\Tinf$ and of $\tuhf$ with its six soft properties
was completed in early July 2026. At that stage the source repository also
contained a completion-tier ``faithfulness'' declaration that, on
inspection, was a placeholder: a named proposition of type
$\texttt{Prop} := \texttt{True}$ discharged by \texttt{trivial}, together
with a theorem whose hypothesis was completion-tier faithfulness restated.
Neither constitutes a proof of anything, and a revision of the companion
ITP note (2026-07-20) explicitly demoted completion-tier faithfulness from
``conditional'' to \emph{open}, documenting the placeholder for
transparency.

The subsequent closing arc, labelled r102--r112 in the repository, then
unfolded in three phases:

\begin{enumerate}[leftmargin=2.4em, itemsep=2pt]
\item \emph{Reduction} (r102--r106): Cauchy--Schwarz for the trace
  pairing, the trace-null two-sided ideal, the Weyl averaging engine,
  algebraic simplicity of the pre-completion tower, and a completion-tier
  averaging operator together reduced faithfulness, in kernel-verified
  form, to a single named hypothesis: \emph{trace-detection of ideals}
  (every nonzero two-sided ideal contains an element of nonzero trace).
\item \emph{Obstruction analysis} (r108): a kernel-verified four-fold
  equivalence showed that trace-detection is not weaker than the goal ---
  it is \emph{logically equivalent} to faithfulness, and to simplicity.
  The reduction had restated the mountain, not descended it. This negative
  result was formalized rather than discarded: it is the precise statement
  of where the classical hard input lives.
\item \emph{Closure} (r109--r112): the trace-preserving conditional
  expectation $E_k$ (the tail partial trace) was constructed on the direct
  limit, proved a $\|\cdot\|_2$-contraction by elementary Cauchy--Schwarz
  (where Kadison--Schwarz was expected to be needed), proved an
  operator-norm contraction via an isometry decomposition (where complete
  positivity was expected to be needed), extended to the completion, and
  used to close faithfulness unconditionally --- with simplicity following
  through the r108 equivalence.
\end{enumerate}

The placeholder was thus first exposed, then reduced, then genuinely
closed. We report this sequence as a strength of the machine-checked
setting: at every stage the kernel-verified statement of \emph{what had
actually been proved} was unambiguous, the gap between claim and theorem
was mechanically visible, and the final theorems are immune to the wishful
reading that produced the original placeholder.

\subsection{A note on naming}\label{sec:naming}

The Lean development is part of a larger source project, and some
identifiers inherit that project's naming scheme: the direct limit is
called \lt{TimelessFieldRing} and its completion
\lt{TimelessFieldCompletion} in the source. These are names, nothing more;
in this paper the objects are denoted $\varinjlim_k A_k$ and $\Tinf$, and
every property claimed of them is one of the kernel-verified theorems
listed below or a standard cited fact. Nothing in this paper depends on
any framework, interpretation, or conjecture of the broader source
project.

\subsection{Outline}

\S\ref{sec:background} recalls the classical background.
\S\ref{sec:tower} constructs the tower, limit, and completion.
\S\ref{sec:trace} constructs $\tuhf$ and its soft properties.
\S\ref{sec:engine} develops the simplicity/faithfulness engine:
Cauchy--Schwarz, the trace-null ideal, Weyl averaging, pre-completion
simplicity, and the four-fold equivalence. \S\ref{sec:condexp} is the
heart: the conditional expectation and the faithfulness proof.
\S\ref{sec:formalization} records formalization notes: mathlib APIs used,
gaps found, and the three contributed API files. \S\ref{sec:reproducibility}
gives the pinned toolchain and verification data. \S\ref{sec:related}
discusses related and future work.

%% =====================================================================
\section{Mathematical background}\label{sec:background}
%% =====================================================================

This section fixes classical terminology for the formalization audience;
all results here are standard~\cite{glimm1960, davidson1996, blackadar2006,
takesaki1979} and are \emph{not} claimed as contributions. Which of them
are formalized, and which are merely cited, is stated precisely in
\S\S\ref{sec:tower}--\ref{sec:condexp}.

\subsection{UHF algebras}

A \emph{uniformly hyperfinite (UHF) algebra} is a unital $C^*$-algebra $A$
containing an increasing sequence of unital $*$-subalgebras
$A_0 \subseteq A_1 \subseteq \cdots$, each $*$-isomorphic to a full matrix
algebra $M_{n_k}(\mathbb{C})$, whose union is dense in $A$. Necessarily
$n_k \mid n_{k+1}$. Glimm~\cite{glimm1960} classified UHF algebras by the
associated supernatural number and proved each is simple with a unique
tracial state; Bratteli~\cite{bratteli1972} extended the framework to
general approximately finite-dimensional (AF) algebras. The algebra
studied here is $M_{3^\infty}$: the case $n_k = 3^k$ with the embeddings
$A \mapsto A \otimes I_3$. The choice of base 3 over the more common CAR
tower $M_{2^\infty}$ is inherited from the source project's surrounding
ternary constructions; as a classification target the two are equally
standard, and every argument below works verbatim in any base $\ge 2$.

\subsection{Tracial states}

A \emph{tracial state} on a unital $C^*$-algebra $A$ is a linear
functional $\tau \colon A \to \mathbb{C}$ with $\tau(1) = 1$,
$\tau(x^*x) \ge 0$, and $\tau(xy) = \tau(yx)$. It is \emph{faithful} if
$\tau(x^*x) = 0$ implies $x = 0$. On $M_n(\mathbb{C})$ the unique tracial
state is the normalised trace $\tau_n = \operatorname{tr}/n$, and its
faithfulness is the statement that
$\operatorname{tr}(M^*M) = \sum_{i,j}|M_{ij}|^2$ vanishes only at
$M = 0$. On a UHF algebra the normalised traces are compatible with the
embeddings and extend to the unique tracial state on the completion; its
faithfulness is classically deduced from simplicity (the closure of the
left kernel is a proper closed two-sided ideal) or from the conditional
expectations onto the finite levels.

\subsection{Simplicity}

A ring is (algebraically) \emph{simple} if its only two-sided ideals are
$0$ and the whole ring. For $C^*$-algebras one often restricts attention
to closed ideals, but UHF algebras satisfy the stronger algebraic
statement, and it is the algebraic statement that is formalized here:
\emph{every} two-sided ideal of $\Tinf$, closed or not, is $\bot$ or
$\top$. The engine of the classical proof is finite-dimensional: averaging
over a suitable finite group of unitaries in $M_n(\mathbb{C})$ projects
any matrix onto its normalised trace times the identity, and since
two-sided ideals absorb conjugation, sums, and scalars, a nonzero ideal
element with nonzero trace forces $1$ into the ideal. The finite unitary
family used here is the discrete Weyl (clock-and-shift)
pair~\cite{weyl1927}: $C = \operatorname{diag}(1, \omega, \dots,
\omega^{n-1})$ with $\omega = e^{2\pi i/n}$, and $S$ the cyclic shift.

\subsection{Conditional expectations}

A \emph{conditional expectation} of a $C^*$-algebra $A$ onto a subalgebra
$B$ is a positive, unital, $B$-bimodule projection $E \colon A \to B$;
by Tomiyama's theorem~\cite{tomiyama1957} these are exactly the norm-one
projections. For a UHF algebra, the canonical conditional expectation onto
the level-$k$ subalgebra is the \emph{partial trace over the tail}: under
$M_{3^{k+1}} \cong M_{3^k} \otimes M_3$, it is
$\mathrm{id} \otimes \tau_3$. Classically, a conditional expectation is
completely positive, hence satisfies the Kadison--Schwarz inequality
$E(x)^*E(x) \le E(x^*x)$~\cite{kadison1952}, from which its
$\|\cdot\|_2$-contractivity with respect to a compatible trace follows,
and its complete positivity together with unitality gives
$\|E\| = 1$ (a fact often organised via the Stinespring
representation~\cite{stinespring1955}). The formalization proves the two
contraction properties it needs \emph{directly}, without any of this
theory; see \S\ref{sec:l2} and \S\ref{sec:opnorm}.

%% =====================================================================
\section{The tower, the limit, and the completion}\label{sec:tower}
%% =====================================================================

\subsection{The tower}

At level $k \in \mathbb{N}$, define
\[
A_k \;:=\; \mathrm{Matrix}~(\mathrm{Fin}~3^k)~(\mathrm{Fin}~3^k)~\mathbb{C},
\]
a finite-dimensional $C^*$-algebra under the $L^2$-operator norm (mathlib's
\lt{Matrix.Norms.L2Operator} scoped instance, which provides the
$C^*$-identity $\|M^*M\| = \|M\|^2$). Successive levels embed by the
ternary tensor construction
\[
\iota_k \colon A_k \hookrightarrow A_{k+1}, \qquad
A \;\longmapsto\; \mathrm{reindex}~(\mathrm{levelStepEquiv}~k)~(A \otimes I_3),
\]
where \lt{levelStepEquiv}~$k$ is the equivalence
$\mathrm{Fin}~3^k \times \mathrm{Fin}~3 \simeq \mathrm{Fin}~3^{k+1}$
induced by the ternary decomposition of $3^{k+1}$, and $\otimes$ is the
Kronecker product. In the Lean source this map is
\lt{substrateEmbedMatrix}, with iterated form \lt{substrateRingHomIter}
$m \to k$ for $m \le k$. Each $\iota_k$ is a unital $*$-ring homomorphism
and an isometry.

\subsection{The inductive limit and its norm}

Mathlib's \lt{Ring.DirectLimit} applied to $\{A_k, \iota_k\}$ produces the
$\mathbb{C}$-algebra
\[
R \;:=\; \varinjlim_k A_k
\qquad (\text{Lean: } \lt{TimelessFieldRing}),
\]
with \texttt{StarRing} structure descending from entrywise conjugate
transpose. Every element of $R$ has a single-level representative: mathlib's
\lt{DirectLimit.exists\_eq\_mk} produces, from $x \in R$, a level $k$ and a
matrix $A \in A_k$ with $x = [\![ (k, A) ]\!]$. Because the
embeddings are isometric, the level norm is well-defined on $R$
(\lt{substrateLevelToTimelessField\_opNorm\_eq}) and makes $R$ a normed
$*$-ring; the canonical level embedding
$\iota_k \colon A_k \to R$ (Lean: \lt{substrateLevelToTimelessField}) is
an isometry (theorem \lt{substrateLevelToTimelessField\_isometry}), and
the cocone identity
$\iota_k \circ (\text{iterated embedding } m \to k) = \iota_m$
holds definitionally up to \texttt{Quotient.sound}
(\lt{substrateLevelToTimelessField\_iter}).

\subsection{The completion}

The metric completion of $R$ under the operator norm is
\[
\Tinf \;:=\; \overline{\varinjlim_k A_k}^{\,\opnorm}
\qquad (\text{Lean: } \lt{TimelessFieldCompletion}),
\]
built with \lt{UniformSpace.Completion}. The algebraic operations extend
by uniform continuity, and the $C^*$-identity passes to the completion by
the classical Cauchy-sequence argument, executed against mathlib's
completion API.

\begin{theorem}[UHF $C^*$-algebra registration]\label{thm:uhf-registration}
$\Tinf$ carries mathlib-native \texttt{Ring}, \texttt{StarRing},
\texttt{NormedRing}, and \texttt{CStarAlgebra} instances, with
$\varinjlim_k A_k$ dense in $\Tinf$. Kernel axiom set:
$[\texttt{propext}, \texttt{Classical.choice}, \texttt{Quot.sound}]$.
\end{theorem}

By construction $\Tinf$ is a UHF algebra in the classical sense: the
norm-closure of an increasing union of full matrix algebras under unital
$*$-embeddings, of supernatural number $3^\infty$. We emphasise that the
mathlib \texttt{CStarAlgebra} instance is the algebra of record; no
appeal is made to the classical classification.

\begin{theorem}[Level-wise simplicity]\label{thm:level-simple}
For every $k$, $A_k = M_{3^k}(\mathbb{C})$ is a mathlib
\texttt{IsSimpleRing} instance, via \lt{DivisionRing.isSimpleRing} on
$\mathbb{C}$ composed with \lt{IsSimpleRing.matrix}.
\end{theorem}

%% =====================================================================
\section{The tracial state and its soft properties}\label{sec:trace}
%% =====================================================================

\subsection{Construction}

On each level the normalised trace
\[
\tau_n \colon A_k \to \mathbb{C}, \qquad
\tau_n(M) := \frac{\operatorname{trace}(M)}{n}, \qquad n = 3^k
\]
is linear, unital, and $1$-Lipschitz for the operator norm (via the
Hilbert--Schmidt bound $\|M\|_{\mathrm{HS}}^2 \le n\|M\|_{\mathsf{op}}^2$
and Cauchy--Schwarz). Compatibility with the embeddings,
$\tau_{3^{k+1}}(\iota_k A) = \tau_{3^k}(A)$, follows from
\lt{Matrix.trace\_kronecker}. The universal property of the direct limit
(\lt{DirectLimit.lift}) produces the pre-trace
$\tau_{\mathrm{pre}} \colon R \to \mathbb{C}$
(Lean: \lt{substrate\_pre\_trace}), additive, unital, and $1$-Lipschitz,
hence uniformly continuous; and
\[
\tuhf \;:=\; \lt{UniformSpace.Completion.extension}\,(\tau_{\mathrm{pre}})
\;\colon\; \Tinf \to \mathbb{C}
\qquad (\text{Lean: } \lt{UHF\_trace})
\]
is its unique continuous extension, agreeing with $\tau_{\mathrm{pre}}$ on
the dense image (\lt{UHF\_trace\_coe}).

\subsection{The soft properties}

\begin{theorem}[The canonical tracial state]\label{thm:tracial-state}
$\tuhf \colon \Tinf \to \mathbb{C}$ satisfies, kernel-verified:
\begin{enumerate}[label=\textup{(\roman*)}, itemsep=1pt]
\item \emph{additive, unital, $1$-Lipschitz, uniformly continuous};
\item \emph{$\mathbb{C}$-linear}: $\tuhf(c \cdot x) = c\,\tuhf(x)$
  (\lt{UHF\_trace\_smul}; bundled as
  \lt{UHF\_trace\_linearMap} $\colon \Tinf \to_{\mathbb{C}} \mathbb{C}$);
\item \emph{tracial}: $\tuhf(xy) = \tuhf(yx)$
  (\lt{UHF\_trace\_mul\_comm});
\item \emph{Hermitian}: $\tuhf(x^*) = \overline{\tuhf(x)}$;
\item \emph{positive}: $0 \le \tuhf(x^*x)$ in the complex order
  (\lt{UHF\_trace\_star\_mul\_self\_nonneg}, under mathlib's scoped
  \texttt{ComplexOrder}).
\end{enumerate}
Kernel axiom set:
$[\texttt{propext}, \texttt{Classical.choice}, \texttt{Quot.sound}]$.
\end{theorem}

Each of (i)--(v) is a closed condition and is proved by the same
three-step pattern: a mathlib matrix lemma at level $k$
(\lt{Matrix.trace\_mul\_comm}, \lt{Matrix.trace\_conjTranspose},
\lt{Matrix.posSemidef\_conjTranspose\_mul\_self} with
\lt{Matrix.PosSemidef.trace\_nonneg}); reduction of direct-limit elements
to single-level representatives via \lt{DirectLimit.exists\_eq\_mk}; and
completion induction (\lt{UniformSpace.Completion.induction\_on}, or its
binary form for (iii)) on the closed set where the identity holds ---
closedness of the positivity set (v) uses
\lt{Complex.orderClosedTopology}.

\subsection{The three tiers of faithfulness}\label{sec:tiers}

Faithfulness stratifies over the construction:

\begin{enumerate}[label=$(\Phi\arabic*)$, leftmargin=3em, itemsep=2pt]
\item \emph{Matrix level.} $\tau_n(M^*M) = 0 \Rightarrow M = 0$, from
  mathlib's \lt{Matrix.trace\_conjTranspose\_mul\_self\_eq\_zero\_iff}
  (Lean: \lt{normalized\_matrix\_trace\_star\_mul\_self\_eq\_zero\_iff}).
  Elementary: the trace of $M^*M$ is the sum of $|M_{ij}|^2$.
\item \emph{Direct limit.}
  $\tau_{\mathrm{pre}}(x^*x) = 0 \Rightarrow x = 0$
  (\lt{substrate\_pre\_trace\_star\_mul\_self\_faithful}): reduce $x$ to a
  single-level representative, apply $(\Phi 1)$, and push $0$ back through
  the embedding.
\item \emph{Completion.} $\tuhf(x^*x) = 0 \Rightarrow x = 0$ for
  $x \in \Tinf$. \textbf{This is the hard tier}, for the reason explained
  in \S\ref{sec:why-hard}: it is not a closed condition, and no density
  argument applies. Its unconditional proof is the subject of
  \S\S\ref{sec:engine}--\ref{sec:condexp}.
\end{enumerate}

%% =====================================================================
\section{The simplicity/faithfulness engine}\label{sec:engine}
%% =====================================================================

This section develops the machinery that first \emph{reduced}
completion-tier faithfulness to a single named hypothesis, and then
\emph{identified} that hypothesis as equivalent to the goal. All results
here are unconditional kernel-verified theorems; the r-labels record the
repository development arc but the content is mathematics, not changelog.

\subsection{Cauchy--Schwarz for the trace pairing (r102)}\label{sec:cs}

Define the trace sesquilinear pairing
$\langle x, y \rangle_\tau := \tuhf(y^* x)$ on $\Tinf$
(Lean: \lt{traceInner}). Positivity (Theorem~\ref{thm:tracial-state}(v))
makes it a positive semi-definite sesquilinear form, and the classical
discriminant argument goes through against mathlib's
\lt{discrim\_le\_zero}:

\begin{theorem}[\lt{UHF\_trace\_cauchySchwarz}]\label{thm:cs}
For all $x, y \in \Tinf$,
\[
\bigl\|\tuhf(y^*x)\bigr\|^2 \;\le\;
\operatorname{Re}\tuhf(x^*x)\cdot \operatorname{Re}\tuhf(y^*y).
\]
\end{theorem}

The proof expands $\tuhf\bigl((x - c y)^*(x - c y)\bigr) \ge 0$ along real
multiples $c = t\,\tuhf(y^*x)$ using completion-tier linearity
(Theorem~\ref{thm:tracial-state}(ii)) and the \texttt{StarModule}
structure, and applies the quadratic discriminant bound. An immediate
corollary is left absorption of the trace-null cone
(\lt{UHF\_trace\_null\_left\_absorption}):
$\tuhf(x^*x) = 0 \Rightarrow \tuhf\bigl((ax)^*(ax)\bigr) = 0$ for every
$a$, since $\tuhf\bigl((ax)^*(ax)\bigr) = \langle a^*ax, x\rangle_\tau$ is
dominated through Theorem~\ref{thm:cs} by
$\operatorname{Re}\tuhf(x^*x) = 0$.

\subsection{The trace-null two-sided ideal (r103)}\label{sec:tracenull}

Let $N := \{x \in \Tinf : \tuhf(x^*x) = 0\}$. Theorem~\ref{thm:cs} makes
$N$ additive (the cross terms in the sesquilinear expansion of
$\tuhf\bigl((x+y)^*(x+y)\bigr)$ are killed by Cauchy--Schwarz), traciality
makes it $*$-closed ($\tuhf(xx^*) = \tuhf(x^*x)$), left absorption plus
$*$-closure gives right absorption, and scalar closure is immediate. The
result is bundled as a mathlib-native object:

\begin{theorem}[\lt{traceNullTwoSidedIdeal}]\label{thm:tracenull}
$N$ is the carrier of a \texttt{TwoSidedIdeal} of $\Tinf$, and it is
proper: $1 \notin N$ since $\tuhf(1) = 1$
\textup{(\lt{traceNullTwoSidedIdeal\_ne\_top})}.
\end{theorem}

This immediately yields the conditional form of faithfulness --- the
honest theorem available at that stage of the development:

\begin{theorem}[\lt{UHF\_trace\_faithful\_of\_simple}]\label{thm:faithful-of-simple}
If every two-sided ideal of $\Tinf$ is $\bot$ or $\top$, then $\tuhf$ is
faithful.
\end{theorem}

\begin{proof}[Proof sketch]
$N$ is a proper two-sided ideal, so simplicity forces $N = \bot$; and
$x \in N$ means exactly $\tuhf(x^*x) = 0$, so membership in $\bot$ gives
$x = 0$. The hypothesis is genuine algebraic simplicity of $\Tinf$ ---
classically Glimm's theorem --- not a self-referential marker.
\end{proof}

\subsection{The Weyl unitary-averaging engine (r104)}\label{sec:weyl}

The engine of every simplicity proof for matrix towers is
finite-dimensional linear algebra. Fix $n \ge 1$ and let
$\omega := e^{2\pi i/n}$ (Lean: \lt{substrateRoot}, a primitive $n$-th
root by \lt{Complex.isPrimitiveRoot\_exp}). Define the \emph{clock} and
\emph{shift} matrices~\cite{weyl1927}
\[
C := \operatorname{diagonal}(j \mapsto \omega^j), \qquad
S_{ij} := [\,i = j + 1\,],
\]
both unitary (\lt{clockMatrix\_mem\_unitaryGroup},
\lt{shiftMatrix\_mem\_unitaryGroup}), with closed-form powers
$C^m = \operatorname{diagonal}(j \mapsto \omega^{jm})$ and
$(S^m)_{ij} = [\,i = j + m\,]$. The character orthogonality relation
$\sum_{a<n} (\omega^i \overline{\omega^j})^a = n\,[\,i = j\,]$
(\lt{substrate\_clock\_character\_sum}, via \lt{geom\_sum\_mul} and
\lt{IsPrimitiveRoot.pow\_inj}) drives a two-stage computation:

\begin{lemma}[Stage A: \lt{clock\_average\_eq\_diagonalPart}]
Conjugation-averaging over the clock powers kills the off-diagonal:
$\frac{1}{n}\sum_{a<n} C^a x\, (C^a)^* =
\operatorname{diagonal}(\operatorname{diag} x)$.
\end{lemma}

\begin{lemma}[Stage B: \lt{shift\_average\_diagonal\_eq\_trace\_smul\_one}]
Conjugation-averaging over the shift powers equalises the diagonal: for
diagonal $D$,
$\frac{1}{n}\sum_{b<n} S^b D\, (S^b)^* =
\bigl(\operatorname{trace} D / n\bigr)\cdot 1$.
\end{lemma}

\begin{theorem}[The engine:
\lt{matrix\_unitary\_average\_eq\_trace\_smul\_one}]\label{thm:weyl}
For every $x \in M_n(\mathbb{C})$,
\[
\frac{1}{n^2}\sum_{a<n}\sum_{b<n}
(C^a S^b)\, x\, (C^a S^b)^* \;=\;
\frac{\operatorname{trace} x}{n}\cdot 1 .
\]
\end{theorem}

Since a set closed under addition, arbitrary left/right multiplication,
and $\mathbb{C}$-scaling absorbs each conjugate $u x u^*$, it absorbs the
average
(\lt{trace\_smul\_one\_mem\_of\_two\_sided\_absorbing}): any two-sided
absorbing set containing $x$ contains
$(\operatorname{trace}x/n)\cdot 1$.

\subsection{Algebraic simplicity below the completion (r105)}\label{sec:presimple}

Firing the engine at a nonzero element requires a nonzero \emph{trace},
which a nonzero ideal element need not have. The fix is the classical
matrix-unit trick: if $x \ne 0$ pick a nonzero entry $x_{pq}$; the matrix
unit $E_{q,p}$ satisfies
$\operatorname{trace}(E_{q,p}\,x) = x_{pq} \ne 0$
(\lt{trace\_single\_mul}), and $E_{q,p}\,x$ is still in the absorbing set.
Averaging and rescaling puts $1$ in the set
(\lt{absorbing\_set\_contains\_one}). Together with
Theorem~\ref{thm:level-simple} this re-derives matrix-level simplicity
in ideal form (\lt{matrix\_twoSidedIdeal\_bot\_or\_top}) and, by pulling a
two-sided ideal $I \le R$ back to each level
(\lt{substrateIdealPullback}, with the four closure lemmas), lifts it to
the direct limit:

\begin{theorem}[\lt{directLimit\_twoSidedIdeal\_bot\_or\_top}]\label{thm:dl-simple}
Every two-sided ideal of $\varinjlim_k M_{3^k}(\mathbb{C})$ is $\bot$ or
$\top$.
\end{theorem}

\subsection{Averaging inside the completion (r106)}\label{sec:avgK}

The Weyl words at level $k$ push forward to unitaries
$u_{a,b} \in \Tinf$, and the completion-tier averaging operator
\[
\operatorname{avg}_k(x) \;:=\; \frac{1}{(3^k)^2}
\sum_{a, b \,<\, 3^k} u_{a,b}\, x\, u_{a,b}^{*}
\qquad (\text{Lean: } \lt{avgK})
\]
is defined for every $x \in \Tinf$. On the image of level $k$ it agrees
exactly with $\tuhf(\cdot)\cdot 1$ by Theorem~\ref{thm:weyl}; combining
density of the finite levels, the non-expansiveness of averaging, and the
$1$-Lipschitz property of $\tuhf$ yields Glimm's averaging estimate
executed inside the completion:

\begin{theorem}[\lt{avgK\_approx}]\label{thm:avgk}
For every $x \in \Tinf$ and $\varepsilon > 0$ there is $k$ with
$\bigl\| \operatorname{avg}_k(x) - \tuhf(x)\cdot 1 \bigr\| < \varepsilon$.
\end{theorem}

Since two-sided ideals absorb each conjugate, $\operatorname{avg}_k$
preserves ideals; if an ideal $I$ contains an element $x$ with
$\tuhf(x) \ne 0$, then Theorem~\ref{thm:avgk} with
$\varepsilon = \|\tuhf(x)\|$ combined with a Neumann-series invertibility
argument forces $1 \in I$. This proves
(\lt{completion\_twoSidedIdeal\_bot\_or\_top\_of\_traceDetects}):
\emph{if} every nonzero two-sided ideal of $\Tinf$ contains an element of
nonzero trace --- the hypothesis named
\lt{SubstrateCompletionTraceDetectsIdeals} --- \emph{then} $\Tinf$ is
simple, hence (Theorem~\ref{thm:faithful-of-simple}) $\tuhf$ is faithful.

\subsection{The four-fold equivalence (r108): locating the mountain}\label{sec:tfae}

Could trace-detection be discharged directly? The formalized answer is a
sharp \emph{no}: it is not a weaker waypoint but a restatement of the
summit.

\begin{theorem}[\lt{traceDetects\_tfae}]\label{thm:tfae}
For the algebra $\Tinf$, the following are equivalent:
\begin{enumerate}[label=\textup{(\arabic*)}, itemsep=1pt]
\item \emph{trace-detection}: every two-sided ideal $I \ne \bot$ contains
  an $x$ with $\tuhf(x) \ne 0$;
\item $\lt{traceNullTwoSidedIdeal} = \bot$;
\item \emph{faithfulness}: $\tuhf(x^*x) = 0 \Rightarrow x = 0$;
\item \emph{simplicity}: every two-sided ideal of $\Tinf$ is $\bot$ or
  $\top$.
\end{enumerate}
The cycle is
$(4) \Rightarrow (3)$ \textup{(Theorem~\ref{thm:faithful-of-simple})},
$(3) \Rightarrow (2)$ \textup{(\lt{traceNull\_eq\_bot\_of\_faithful})},
$(2) \Rightarrow (1)$ \textup{(\lt{traceDetects\_of\_traceNull\_eq\_bot})},
$(1) \Rightarrow (4)$ \textup{(\S\ref{sec:avgK})}, formalized with
mathlib's \texttt{List.TFAE}.
\end{theorem}

Theorem~\ref{thm:tfae} is the honest obstruction map: no hypothesis
genuinely weaker than faithfulness itself can close the loop, and the
classical hard input --- Glimm's theorem, or something of equal strength
--- is really required. The same development also probed the classical
spectral-projection route (producing, from a nonzero selfadjoint element
of an ideal, a nonzero projection in the ideal via the continuous
functional calculus) and formalized exactly where it stalls in the
non-unital cfc setting; the reusable API built for that probe is described
in \S\ref{sec:formathlib}. The route that succeeded is different, and is
the subject of the next section.

%% =====================================================================
\section{The conditional expectation and the faithfulness proof}\label{sec:condexp}
%% =====================================================================

The closing argument constructs the canonical conditional expectation
onto each finite level, proves two contraction properties for it by
elementary means, extends it to the completion, and derives faithfulness
from the reconstruction property $E_k x \to x$. Simplicity then follows
through Theorem~\ref{thm:tfae}. Throughout, $n = 3^k$.

\subsection[The partial trace E\_k (r109)]{The partial trace $E_k$ (r109)}\label{sec:Ek}

Under the identification
$M_{3^{n+1}}(\mathbb{C}) \cong M_{3^n}(\mathbb{C}) \otimes M_3(\mathbb{C})$
furnished by \lt{levelStepEquiv}, the single-step partial trace
$\mathrm{id} \otimes \tau_3$ is, entrywise,
\[
(\mathcal{E} B)_{ij} \;=\; \frac{1}{3}\sum_{t=0}^{2}
B_{(i,t),(j,t)}
\qquad (\text{Lean: } \lt{partialTraceStep}),
\]
where $(i,t)$ abbreviates $\mathrm{levelStepEquiv}\,(i,t)$. Iterating
$j - k$ steps gives
$\lt{partialTraceDown} \colon M_{3^j} \to M_{3^k}$ for $k \le j$, and
descending through single-level representatives defines
\[
E_k \colon \varinjlim_m A_m \longrightarrow M_{3^k}(\mathbb{C})
\qquad (\text{Lean: } \lt{condExp}).
\]
The kernel-verified properties are the trace-preserving retraction
properties needed for the faithfulness argument below; we call $E_k$ the
\emph{level-$k$ expectation map}. We emphasize that the full abstract
conditional-expectation API --- positivity, the $A_k$-bimodule identity
$E_k(a x b) = a\,E_k(x)\,b$ for $a,b \in A_k$, and idempotence as a
norm-one projection --- is \emph{not} formalized here; only the
properties listed in Theorem~\ref{thm:Ek} are, and only those are used.
The name ``conditional expectation'' below is thus informal shorthand for
this partial-trace retraction, not a claim that the classical
conditional-expectation axioms are machine-checked:

\begin{theorem}[Properties of $E_k$]\label{thm:Ek}
$E_k$ is additive and $\mathbb{C}$-homogeneous, unital, star-preserving,
inverts the embedding on the level-$k$ corner
\textup{(\lt{condExp\_level\_k}: $E_k(\iota_k A) = A$, and more generally
\lt{condExp\_lower\_level} for levels $m \le k$)}, and is
trace-preserving: $\tau_{3^k}(E_k x) = \tau_{\mathrm{pre}}(x)$
\textup{(\lt{condExp\_normalized\_trace})}.
\end{theorem}

\subsection[The 2-norm contraction, elementarily (r110)]{The $\|\cdot\|_2$-contraction, elementarily (r110)}\label{sec:l2}

The first analytic requirement is the contraction of $E_k$ in the trace
$2$-seminorm:
\begin{equation}\label{eq:l2}
\operatorname{Re}\tau_{3^k}\bigl((E_k x)^*(E_k x)\bigr)
\;\le\;
\operatorname{Re}\tau_{\mathrm{pre}}(x^*x)
\qquad (\text{Lean: } \lt{condExp\_l2\_contraction}).
\end{equation}
Classically \eqref{eq:l2} is the Kadison--Schwarz
inequality~\cite{kadison1952} for the completely positive map $E_k$,
composed with the trace. Mathlib has neither Kadison--Schwarz nor any
complete-positivity theory, and formalizing them would have been a
project in itself. The formalization instead targets the \emph{scalar}
inequality directly, and it turns out to be elementary. For the single
step: by finite Cauchy--Schwarz over three terms,
\[
\Bigl|\sum_{t=0}^{2} a_t\Bigr|^2 \le 3 \sum_{t=0}^{2}|a_t|^2
\qquad (\lt{normSq\_sum\_fin3\_le}),
\]
so each entry of $\mathcal{E}B$ obeys
$|(\mathcal{E}B)_{ij}|^2 \le \tfrac{1}{3}\sum_t |B_{(i,t),(j,t)}|^2$.
Summing over $i, j$, the right-hand side is precisely the
\emph{diagonal-tail sub-sum} ($s = t$) of
\[
\tau_{3^{n+1}}(B^*B) \;=\; \frac{1}{3^{n+1}}
\sum_{i,s,j,t} \bigl|B_{(j,t),(i,s)}\bigr|^2,
\]
all of whose terms are nonnegative; mathlib's
\lt{Finset.single\_le\_sum} (in its sub-sum form) closes the single-step
inequality
$\tau_{3^n}\bigl((\mathcal{E}B)^*(\mathcal{E}B)\bigr) \le
\tau_{3^{n+1}}(B^*B)$
(\lt{partialTraceStep\_normSq\_trace\_le}). Iterating through
\lt{partialTraceDown} and descending to the direct limit gives
\eqref{eq:l2}. No positive-semidefinite operator theory and no complete
positivity enter at any point --- a useful proof-engineering
simplification: for trace-compatible partial traces, the scalar
$2$-norm contraction is a three-line Cauchy--Schwarz fact, so the
Kadison--Schwarz detour that the textbook proof takes is avoidable in
the formalization.

\subsection{The operator-norm contraction via isometries (r111)}\label{sec:opnorm}

Extending $E_k$ to the completion requires uniform continuity, i.e.\ the
operator-norm bound $\|E_k x\| \le \|x\|$. Classically this is
``a unital completely positive map has norm one,'' again unavailable. The
formalized proof uses the isometry (Stinespring-flavoured
\cite{stinespring1955}) decomposition of the partial trace. For
$t \in \{0,1,2\}$ let
$W_t \colon \mathbb{C}^{3^n} \to \mathbb{C}^{3^{n+1}}$ be the coordinate
injection $v \mapsto v \otimes e_t$, as a
$3^{n+1} \times 3^n$ matrix (Lean: \lt{Winj}). Then
\begin{equation}\label{eq:stinespring}
\mathcal{E} B \;=\; \frac{1}{3}\sum_{t=0}^{2} W_t^{*}\, B\, W_t
\qquad (\lt{partialTraceStep\_eq\_smul\_sum}),
\end{equation}
and each $W_t$ satisfies $W_t^* W_t = I$
(\lt{Winj\_conjTranspose\_mul\_self}). The $C^*$-identity for the
$L^2$-operator norm gives
$\|W_t\|^2 = \|W_t^*W_t\| = \|I\| \le 1$
(\lt{Winj\_opNorm\_sq\_le\_one}, using
\lt{Matrix.l2\_opNorm\_conjTranspose\_mul\_self} and
\lt{ContinuousLinearMap.norm\_id\_le}), so submultiplicativity and
$\|W_t^*\| = \|W_t\|$ yield
$\|W_t^* B W_t\| \le \|B\|$, and averaging over the three slots closes
\[
\|\mathcal{E}B\| \le \|B\|
\qquad (\lt{partialTraceStep\_opNorm\_le}),
\]
hence, iterating, $\|E_k x\| \le \|x\|$ (\lt{condExp\_opNorm\_le}).
Together with additivity this makes $E_k$ $1$-Lipschitz
(\lt{condExp\_lipschitz}), hence uniformly continuous. The only mathlib
inputs are four $L^2$-operator-norm lemmas and the norm bound on the
identity operator; no completely-positive-map theory is used.

\subsection{Extension to the completion and the closing argument (r112)}\label{sec:summit}

Uniform continuity extends $E_k$ to the completion in two guises. First,
the matrix-valued extension
\[
\lt{condExpFin}\;k \colon \Tinf \to M_{3^k}(\mathbb{C}),
\qquad \lt{condExpFin}\,k := \lt{UniformSpace.Completion.extension}\,(E_k).
\]
Second, its recoordinatization back into the completion,
\[
\lt{condExpCompletion}\;k \colon \Tinf \to \Tinf,
\qquad x \mapsto \iota_k\bigl(\lt{condExpFin}\,k\,x\bigr).
\]
Both are $1$-Lipschitz (\lt{condExpFin\_lipschitz},
\lt{condExpCompletion\_lipschitz}, the latter through the isometry of
$\iota_k$, Theorem in \S\ref{sec:tower}), and agreeing with $E_k$ on the
dense image (\lt{condExpFin\_coe}, \lt{condExpCompletion\_coe}). Three
statements assemble the proof.

\begin{lemma}[Reconstruction: \lt{condExpCompletion\_tendsto}]\label{lem:tendsto}
For every $x \in \Tinf$,
$\lt{condExpCompletion}\;k\;x \longrightarrow x$ as $k \to \infty$.
\end{lemma}

\begin{proof}[Proof sketch]
An $\varepsilon/2$ density lift over the uniformly $1$-Lipschitz family:
approximate $x$ within $\varepsilon/2$ by a dense-image element with
single-level representative at level $m$; for $k \ge m$ the map
$\lt{condExpCompletion}\;k$ \emph{fixes} that element
(\lt{condExpCompletion\_fixes\_coe\_lower}, from
\lt{condExp\_lower\_level} and the cocone identity), and the Lipschitz
bound controls the difference at $x$.
\end{proof}

\begin{lemma}[The lifted $\|\cdot\|_2$-contraction:
\lt{condExpCompletion\_l2\_contraction}]\label{lem:l2-lift}
For every $k$ and $x \in \Tinf$,
\[
\operatorname{Re}\tuhf\bigl((\lt{condExpCompletion}\,k\,x)^*
(\lt{condExpCompletion}\,k\,x)\bigr)
\;\le\;
\operatorname{Re}\tuhf(x^*x).
\]
\end{lemma}

\begin{proof}[Proof sketch]
The inequality between two continuous real-valued functions of $x$ holds
on the dense image by \eqref{eq:l2} and the trace bridge
$\tuhf\bigl((E_kx)^*(E_kx)\bigr) =
\tau_{3^k}\bigl((\lt{condExpFin}\,k\,x)^*(\lt{condExpFin}\,k\,x)\bigr)$
(\lt{UHF\_trace\_condExpCompletion}, definitional since
\lt{condExpCompletion} factors through $\iota_k$); completion induction on
the closed set $\{x : f(x) \le g(x)\}$ (\lt{isClosed\_le}) transfers it to
all of $\Tinf$.
\end{proof}

\begin{lemma}[Matrix faithfulness in real-part form:
\lt{matrix\_faithful\_of\_re\_nonpos}]\label{lem:re-form}
If $\operatorname{Re}\tau_n(M^*M) \le 0$ then $M = 0$.
\end{lemma}

\begin{proof}[Proof sketch]
Positivity pins $\tau_n(M^*M)$ to the nonnegative reals; a nonpositive
real part then forces $\tau_n(M^*M) = 0$, and matrix-tier faithfulness
$(\Phi 1)$ concludes.
\end{proof}

\begin{theorem}[The summit: \lt{UHF\_trace\_faithful}]\label{thm:faithful}
For every $x \in \Tinf$: if $\tuhf(x^*x) = 0$ then $x = 0$. Kernel axiom
set:
$[\texttt{propext}, \texttt{Classical.choice}, \texttt{Quot.sound}]$;
zero project axioms.
\end{theorem}

\begin{proof}
Suppose $\tuhf(x^*x) = 0$. By Lemma~\ref{lem:l2-lift} and the trace
bridge, $\operatorname{Re}\tau_{3^k}\bigl((\lt{condExpFin}\,k\,x)^*
(\lt{condExpFin}\,k\,x)\bigr) \le 0$ for every $k$; by
Lemma~\ref{lem:re-form}, $\lt{condExpFin}\,k\,x = 0$, hence
$\lt{condExpCompletion}\,k\,x = 0$, for every $k$. By
Lemma~\ref{lem:tendsto} the constant-zero sequence
$\lt{condExpCompletion}\,k\,x$ converges to $x$; uniqueness of limits
(\lt{tendsto\_nhds\_unique}) gives $x = 0$.
\end{proof}

\begin{corollary}[Glimm simplicity, machine-checked:
\lt{substrate\_completion\_simple\_unconditional}]\label{cor:simple}
Every two-sided ideal of $\Tinf$ is $\bot$ or $\top$.
\end{corollary}

\begin{proof}
Theorem~\ref{thm:faithful} is item (3) of the four-fold equivalence
(Theorem~\ref{thm:tfae}); item (4) follows. In Lean the step is literally
\lt{traceDetects\_iff\_simple.mp}
$\circ$ \lt{traceDetects\_iff\_faithful.mpr} applied to
\lt{UHF\_trace\_faithful}.
\end{proof}

\begin{remark}
The classical order of deduction is inverted: Glimm proves simplicity
first and derives faithfulness of the trace; here faithfulness is proved
directly through the conditional expectations, and simplicity follows.
Both directions of the equivalence are formalized, so the classical order
is also machine-checked --- Theorem~\ref{thm:faithful-of-simple} is
exactly Glimm's implication. Note also that Corollary~\ref{cor:simple} is
\emph{algebraic} simplicity (all two-sided ideals, not only closed ones),
the stronger classical statement, which holds for all UHF
algebras~\cite{davidson1996}.
\end{remark}

\begin{remark}[Residuals discharged]
Theorem~\ref{thm:faithful} discharges, as kernel theorems, every named
hypothesis left standing earlier in the arc: the conditional-expectation
extension package
(\lt{condExpCompletionExtension\_holds}), the r100 completion-faithfulness
residual
(\lt{substrate\_UHF\_completion\_positive\_faithfulness}), and
trace-detection
(\lt{substrate\_completion\_traceDetectsIdeals}). The capstone theorem
\lt{r112\_substrate\_completion\_faithful\_capstone} bundles the seven
statements --- extension, reconstruction, lifted contraction,
faithfulness, the two residuals, and simplicity --- as a single verified
conjunction.
\end{remark}

\subsection{Uniqueness of the trace (r113)}\label{sec:uniqueness}

The same conditional-expectation machinery yields the third leg of the
classical UHF identification: $\tau_{\mathsf{UHF}}$ is the \emph{only}
tracial state on $\Tinf$. We package the hypotheses as a predicate
\lt{IsTracialState}\,$\varphi$ requiring $\varphi$ to be continuous,
additive, $\mathbb{C}$-homogeneous, unital ($\varphi(1)=1$), and tracial
($\varphi(xy)=\varphi(yx)$); positivity is deliberately omitted.

\begin{theorem}[Uniqueness of the substrate trace]\label{thm:unique}
For every $\varphi$ with \lt{IsTracialState}\,$\varphi$ and every
$x \in \Tinf$, $\varphi(x) = \tau_{\mathsf{UHF}}(x)$
\textup{(\lt{substrate\_UHF\_trace\_unique})}. Kernel axioms
$[\texttt{propext},\texttt{Classical.choice},\texttt{Quot.sound}]$.
\end{theorem}

The proof is the standard argument, formalized in three steps. \emph{Finite
level:} any additive, $\mathbb{C}$-homogeneous, unital, tracial functional
on $M_n(\mathbb{C})$ equals the normalized trace
(\lt{matrix\_tracial\_state\_unique}) --- the off-diagonal matrix units are
commutators $E_{ij} = [E_{ik},E_{kj}]$ and hence annihilated, while the
diagonal units are cyclic-conjugate and hence equal, so $\varphi$ is pinned
to $\varphi(1)/n$ times the trace; unitality fixes the constant. Mathlib
does not carry this lemma, so it is proved directly from
\lt{Matrix.single} identities and \lt{matrix\_eq\_sum\_single}. \emph{Level
lift:} for $x = \iota_k(A)$, $\varphi \circ \iota_k$ is such a functional,
so $\varphi(x) = \tau_{3^k}(A) = \tau_{\mathsf{UHF}}(x)$. \emph{Density
lift:} $\varphi$ and $\tau_{\mathsf{UHF}}$ are continuous and agree on the
dense direct-limit image, so they agree everywhere
(\lt{isClosed\_eq} on the completion, exactly as in
\S\ref{sec:summit}).

Combining Theorem~\ref{thm:faithful}, the simplicity corollary, and
Theorem~\ref{thm:unique}, the capstone
\lt{r113\_substrate\_UHF\_factor\_capstone} records that $\Tinf$ is a
faithful, simple, uniquely-traced UHF $C^*$-algebra --- i.e.\ the Glimm
$3^\infty$ UHF factor.

%% =====================================================================
\section{Formalization notes}\label{sec:formalization}
%% =====================================================================

\subsection{Mathlib infrastructure used}

The construction is deliberately mathlib-native; the load-bearing APIs
were:

\begin{itemize}[itemsep=2pt]
\item \lt{Ring.DirectLimit} with \lt{DirectLimit.lift} and
  \lt{DirectLimit.exists\_eq\_mk} (single-level representatives, used in
  essentially every direct-limit argument);
\item \lt{UniformSpace.Completion}: \lt{extension},
  \lt{extension\_coe}, \lt{induction\_on} and \lt{induction\_on}$_2$
  (closed-set completion induction), \lt{denseRange\_coe},
  \lt{coe\_isometry};
\item the scoped $L^2$-operator-norm $C^*$-structure on matrices
  (\lt{Matrix.Norms.L2Operator}): \lt{l2\_opNorm\_mul},
  \lt{l2\_opNorm\_conjTranspose},
  \lt{l2\_opNorm\_conjTranspose\_mul\_self}, \lt{cstar\_norm\_def};
\item matrix trace and positivity lemmas: \lt{Matrix.trace\_kronecker},
  \lt{trace\_mul\_comm}, \lt{trace\_conjTranspose},
  \lt{trace\_diagonal},
  \lt{trace\_conjTranspose\_mul\_self\_eq\_zero\_iff},
  \lt{posSemidef\_conjTranspose\_mul\_self},
  \lt{PosSemidef.trace\_nonneg};
\item \lt{IsSimpleRing.matrix} and \lt{TwoSidedIdeal} (lattice
  structure, \lt{mk'}, \lt{eq\_bot\_or\_eq\_top});
\item roots of unity (\lt{Complex.isPrimitiveRoot\_exp},
  \lt{IsPrimitiveRoot.pow\_inj}, \lt{geom\_sum\_mul}) for the Weyl
  engine;
\item the scoped \texttt{ComplexOrder} with
  \lt{Complex.orderClosedTopology} (positivity as a closed condition),
  and \lt{discrim\_le\_zero} (Cauchy--Schwarz);
\item for the API files of \S\ref{sec:formathlib}: the non-unital
  continuous functional calculus (\lt{cfc}$_{\mathrm{n}}$,
  \lt{cfc}$_{\mathrm{n}}$\lt{Hom}), \lt{CStarAlgebra.approximateUnit}, and the
  Stone--Weierstrass induction principle
  \lt{ContinuousMapZero.induction\_on\_of\_compact}.
\end{itemize}

\subsection{Gaps found in mathlib}\label{sec:gaps}

Three absences shaped the mathematics of
\S\S\ref{sec:engine}--\ref{sec:condexp} and are, we suggest, the natural
next targets for mathlib's operator-algebra development:

\begin{enumerate}[label=\textbf{(G\arabic*)}, leftmargin=3em, itemsep=3pt]
\item \textbf{No completely positive map theory.} Mathlib has no notion
  of complete positivity, no Kadison--Schwarz inequality, no partial
  trace, and no ``unital CP map has norm one.'' This forced --- and, in
  the end, rewarded --- the elementary routes of \S\ref{sec:l2}
  (three-term Cauchy--Schwarz instead of Kadison--Schwarz) and
  \S\ref{sec:opnorm} (isometry decomposition instead of CP norm theory).
  The two proofs are shorter than the classical ones and need a fraction
  of the theory; they may be of expository interest independent of the
  formalization.
\item \textbf{No spectral projections.} The Borel functional calculus and
  spectral projections are absent; the continuous functional calculus
  exists but produces projections only from \emph{clopen} spectral sets.
  The Route-A probe of \S\ref{sec:tfae} built exactly that
  (\S\ref{sec:formathlib}) and identified finite-quasispectrum witnesses
  as the missing hypothesis --- a wall that the conditional-expectation
  route then bypassed entirely.
\item \textbf{No $C^*$-order on completions.} $\Tinf$, built as a
  \lt{UniformSpace.Completion}, does not automatically carry the
  \texttt{PartialOrder}/\texttt{StarOrderedRing} structure of a
  $C^*$-algebra order; positivity of $\tuhf$ was therefore formulated and
  proved scalar-side, in the \texttt{ComplexOrder} on $\mathbb{C}$,
  rather than operator-side. An order-instance transfer for completions
  of normed $*$-algebras would let statements like
  Theorem~\ref{thm:tracial-state}(v) take their textbook form.
\end{enumerate}

\subsection{Contributed API: the \texttt{ForMathlib} files}\label{sec:formathlib}

Three standalone files, written to mathlib contribution standards
(mathlib copyright header, full docstrings, no project dependencies
beyond one another), were produced during the Route-A probe and are
offered upstream:

\begin{itemize}[itemsep=3pt]
\item \lt{PF/ForMathlib/TwoSidedIdealClosure.lean} ---
  \lt{TwoSidedIdeal.closure}: the topological closure of a two-sided
  ideal in a (non-unital, non-associative) topological ring is a
  two-sided ideal, with the expected API (\lt{le\_closure},
  \lt{closure\_minimal}, \lt{closure\_eq\_of\_isClosed}) and, in a unital
  topological ring with open unit group (e.g.\ a Banach algebra),
  properness of the closure of a proper ideal (\lt{closure\_ne\_top}).
  Mirrors mathlib's \lt{Ideal.closure} for one-sided ideals.
\item \lt{PF/ForMathlib/CfcMemTwoSidedIdeal.lean} --- the non-unital
  continuous functional calculus maps into closed two-sided ideals:
  for $a$ selfadjoint in a closed two-sided ideal $I$ of a non-unital
  $C^*$-algebra and \emph{any} $f \colon \mathbb{R} \to \mathbb{R}$,
  $\lt{cfc}_{\mathrm{n}}\,f\,a \in I$
  (\lt{TwoSidedIdeal.cfc}$_{\mathrm{n}}$\lt{\_mem\_of\_isClosed}), unconditional thanks to
  the junk-value convention. The proof runs an approximate-unit argument
  for scalar closure (\lt{smul\_mem\_of\_isClosed}) and Stone--Weierstrass
  induction over $C(\sigma_{\mathrm{n}}\,\mathbb{R}\,a, \mathbb{R})_0$.
\item \lt{PF/ForMathlib/ClopenSpectralProjection.lean} ---
  \lt{spectralProjection}~$a$~$U := \lt{cfc}_{\mathrm{n}}
  (U.\mathrm{indicator}\,1)\,a$ for $U$ clopen in the quasispectrum with
  $0 \notin U$: selfadjoint, idempotent, nonzero iff $U$ meets the
  quasispectrum, and a member of every closed two-sided ideal containing
  $a$ --- i.e., clopen spectral data away from $0$ manufactures a genuine
  projection inside a nonzero ideal.
\end{itemize}

\subsection{Companion result}

The same repository formalizes, independently of everything above, a
sesquilinear self-adjointness identity for a symmetrised base-3 transfer
operator on the mathlib Hilbert space $L^2([0,1], dx/x)$
(theorem \lt{T3\_self\_adjoint\_conj}; see the companion ITP note and
\cite{cohen2025transferOp}). It shares the base-3 architecture but is
logically disjoint from the UHF development: neither result uses the
other. We record it here only for completeness of the repository's
verified operator-algebra content.

%% =====================================================================
\section{Reproducibility}\label{sec:reproducibility}
%% =====================================================================

\paragraph{Pinned toolchain.}
\begin{description}[font=\normalfont\bfseries,leftmargin=2em,itemsep=2pt]
\item[Lean toolchain] \texttt{leanprover/lean4:v4.24.0-rc1}
\item[mathlib revision] commit \texttt{eed770a434957369c6262aa3fb1d6426419016d4}
\item[Repository] \href{https://github.com/FractalDevTeam/Principia-Fractalis}{\texttt{github.com/FractalDevTeam/Principia-Fractalis}}, branch \texttt{master}
\item[HEAD reference (verified)] \texttt{91cc1a01} (2026-07-23); build and \texttt{\#print axioms} run same-day
\end{description}

\paragraph{Build commands.}
\begin{verbatim}
git clone https://github.com/FractalDevTeam/Principia-Fractalis.git
cd Principia-Fractalis/PF_Lean4_Code
elan install $(cat lean-toolchain)
lake build PF     # 4,472 jobs
lake build        # 8,930 jobs (adds PF_Lean4Lean re-elaboration)
\end{verbatim}

\paragraph{Headline theorems and \texttt{\#print axioms} output.}
The axiom audit blocks at the end of
\lt{PF/SubstrateCompletionFaithful.lean} (and the corresponding blocks in
the r102--r111 files) report, for every theorem cited in this paper:

\begin{small}
\begin{verbatim}
'PrincipiaTractalis.SubstrateCompletionFaithful.UHF_trace_faithful'
  depends on axioms: [propext, Classical.choice, Quot.sound]

'PrincipiaTractalis.SubstrateCompletionFaithful.
  substrate_completion_simple_unconditional'
  depends on axioms: [propext, Classical.choice, Quot.sound]

'PrincipiaTractalis.SubstrateCompletionFaithful.
  r112_substrate_completion_faithful_capstone'
  depends on axioms: [propext, Classical.choice, Quot.sound]
\end{verbatim}
\end{small}

These are the three standard Lean~4 kernel axioms (propositional
extensionality, choice, quotient soundness); there are no project-level
axioms, and zero \texttt{sorry} occurrences on any code path leading to
the theorems above. The \lt{PF\_Lean4Lean} package (separate
\texttt{lakefile.toml} and package hash, same pinned mathlib)
re-elaborates each theorem through the production Lean~4 kernel under a
distinct package boundary, adding roughly $4{,}400$ further verification
jobs.

\paragraph{Where to find each result.}
Table~\ref{tab:files} maps the paper's sections to source files.

\begin{table}[ht]
\centering
\small
\begin{tabular}{@{}lll@{}}
\toprule
Content & Paper & Lean file (under \texttt{PF/}) \\
\midrule
Tower, limit, completion, instances & \S\ref{sec:tower} & \lt{Substrate*} (r41--r60 chain) \\
Tracial state and soft properties & \S\ref{sec:trace} & \lt{SubstrateUHFTraceIs*.lean} \\
Cauchy--Schwarz for the trace & \S\ref{sec:cs} & \lt{SubstrateUHFTraceCauchySchwarz.lean} \\
Trace-null ideal; faithful-of-simple & \S\ref{sec:tracenull} & \lt{SubstrateUHFTraceNullIdeal.lean} \\
Weyl averaging engine & \S\ref{sec:weyl} & \lt{SubstrateMatrixUnitaryAveraging.lean} \\
Direct-limit simplicity & \S\ref{sec:presimple} & \lt{SubstrateDirectLimitSimplicity.lean} \\
Completion averaging & \S\ref{sec:avgK} & \lt{SubstrateCompletionSimplicity.lean} \\
Four-fold equivalence & \S\ref{sec:tfae} & \lt{SubstrateTraceDetectionAttempt.lean} \\
Conditional expectation $E_k$ & \S\ref{sec:Ek} & \lt{SubstrateConditionalExpectation.lean} \\
$\|\cdot\|_2$-contraction & \S\ref{sec:l2} & \lt{SubstrateCondExpContraction.lean} \\
Operator-norm contraction & \S\ref{sec:opnorm} & \lt{SubstrateCondExpOpNorm.lean} \\
Extension; faithfulness; simplicity & \S\ref{sec:summit} & \lt{SubstrateCompletionFaithful.lean} \\
Contributed API & \S\ref{sec:formathlib} & \lt{ForMathlib/*.lean} \\
\bottomrule
\end{tabular}
\medskip
\caption{Section-to-source map.}\label{tab:files}
\end{table}

%% =====================================================================
\section{Related and future work}\label{sec:related}
%% =====================================================================

\subsection{Related work}

Mathlib's $C^*$-algebra hierarchy, continuous functional calculus, matrix
positivity infrastructure, direct limits, and completions --- all
community developments, with the functional-analysis foundations described
in~\cite{dupuis-lewis-macbeth2022} --- are the platform this work builds
on and would have been impossible without. To the best of our knowledge,
however --- based on keyword and definition searches of the mathlib
library, the Coq \texttt{opam} package archive, and the Isabelle Archive
of Formal Proofs for UHF/AF/Glimm/simple-$C^*$/tracial-state content, and
on the operator-algebra formalization literature --- no prior
formalization in Lean, Coq, or Isabelle contains a concrete
infinite-dimensional simple $C^*$-algebra, a canonical faithful tracial
state on one, or any version of Glimm's simplicity
theorem~\cite{glimm1960}. This does not constitute an exhaustive survey of
every ITP library, and we would welcome correction from operator-algebra
formalization specialists. On the classical
side, everything proved here is textbook
material~\cite{davidson1996, blackadar2006, takesaki1979}: Glimm's
classification and simplicity~\cite{glimm1960}, Powers' work on the
associated factors~\cite{powers1967}, Bratteli's AF
framework~\cite{bratteli1972}, and the conditional-expectation theory of
Tomiyama~\cite{tomiyama1957} descending from
Kadison~\cite{kadison1952} and Stinespring~\cite{stinespring1955}. The
contribution is not new mathematics about $M_{3^\infty}$; it is the first
machine-checked path to these facts, together with the elementary
replacement proofs (\S\ref{sec:l2}, \S\ref{sec:opnorm}) that the
machine-checked setting demanded and produced.

\subsection{Future work}

Natural next steps, roughly in order of proximity:

\begin{itemize}[itemsep=2pt]
\item \emph{Upstreaming}: the three \texttt{ForMathlib} files
  (\S\ref{sec:formathlib}), and de-specialising the tower construction
  from base 3 to an arbitrary supernatural number --- most proofs are
  already base-generic in structure.
\item \emph{Uniqueness of the trace}: $\tuhf$ is canonical by
  construction; formalizing that it is the \emph{only} tracial state on
  $\Tinf$ (classically via the averaging engine of \S\ref{sec:weyl},
  which is already formalized) would complete the textbook triple
  simple/unique-trace/faithful.
\item \emph{CP-map foundations}: gap (G1) --- partial traces, complete
  positivity, Kadison--Schwarz --- now has a worked concrete example to
  test an abstract API against.
\item \emph{GNS and the hyperfinite II$_1$ factor}: the GNS
  representation of $\tuhf$, the $\|\cdot\|_2$-completion, and the trace
  on the double commutant, connecting to the Murray--von Neumann
  hyperfinite factor~\cite{murray-vonneumann1936, powers1967}.
\item \emph{Glimm's classification}: isomorphism invariance via the
  supernatural number, for which the clopen spectral projection API of
  \S\ref{sec:formathlib} is a first ingredient.
\end{itemize}

%% =====================================================================
\section*{Acknowledgments}
%% =====================================================================

The Lean~4 and mathlib community built the $C^*$-algebra, direct-limit,
completion, and functional-calculus infrastructure on which every theorem
in this paper stands.

\begin{thebibliography}{99}

\bibitem{blackadar2006}
B.~Blackadar.
\newblock \emph{Operator Algebras: Theory of $C^*$-Algebras and von Neumann Algebras}.
\newblock Encyclopaedia of Mathematical Sciences 122, Springer, 2006.

\bibitem{bratteli1972}
O.~Bratteli.
\newblock Inductive limits of finite dimensional $C^*$-algebras.
\newblock \emph{Trans. Amer. Math. Soc.}, 171:195--234, 1972.

\bibitem{cohen2025transferOp}
P.~Cohen.
\newblock A Modified Transfer Operator Approach to the Riemann Hypothesis: Numerical Evidence and Theoretical Framework.
\newblock Manuscript, June 12, 2025. Preserved at \texttt{Papers/PriorWork\_Cohen2025\_TransferOperatorRH/} in the source repository.

\bibitem{davidson1996}
K.~R.~Davidson.
\newblock \emph{$C^*$-Algebras by Example}.
\newblock Fields Institute Monographs 6, American Mathematical Society, 1996.

\bibitem{dupuis-lewis-macbeth2022}
F.~Dupuis, R.~Y.~Lewis, and H.~Macbeth.
\newblock Formalized functional analysis with semilinear maps.
\newblock In \emph{13th International Conference on Interactive Theorem Proving (ITP 2022)}, LIPIcs 237, 2022.

\bibitem{glimm1960}
J.~Glimm.
\newblock On a certain class of operator algebras.
\newblock \emph{Trans. Amer. Math. Soc.}, 95:318--340, 1960.

\bibitem{kadison1952}
R.~V.~Kadison.
\newblock A generalized Schwarz inequality and algebraic invariants for operator algebras.
\newblock \emph{Ann. of Math.~(2)}, 56:494--503, 1952.

\bibitem{mathlib2020}
The mathlib Community.
\newblock The Lean Mathematical Library.
\newblock In \emph{Proc.\ CPP 2020}, pages 367--381, 2020.

\bibitem{moura-ullrich2021}
L.~de Moura and S.~Ullrich.
\newblock The Lean~4 Theorem Prover and Programming Language.
\newblock In \emph{CADE 28}, 2021.

\bibitem{murray-vonneumann1936}
F.~J.~Murray and J.~von Neumann.
\newblock On rings of operators.
\newblock \emph{Ann. of Math.~(2)}, 37(1):116--229, 1936.

\bibitem{powers1967}
R.~T.~Powers.
\newblock Representations of uniformly hyperfinite algebras and their associated von Neumann rings.
\newblock \emph{Ann. of Math.~(2)}, 86:138--171, 1967.

\bibitem{stinespring1955}
W.~F.~Stinespring.
\newblock Positive functions on $C^*$-algebras.
\newblock \emph{Proc. Amer. Math. Soc.}, 6:211--216, 1955.

\bibitem{takesaki1979}
M.~Takesaki.
\newblock \emph{Theory of Operator Algebras I}.
\newblock Springer, 1979.

\bibitem{tomiyama1957}
J.~Tomiyama.
\newblock On the projection of norm one in $W^*$-algebras.
\newblock \emph{Proc. Japan Acad.}, 33:608--612, 1957.

\bibitem{weyl1927}
H.~Weyl.
\newblock Quantenmechanik und Gruppentheorie.
\newblock \emph{Zeitschrift f\"ur Physik}, 46:1--46, 1927.

\end{thebibliography}

\end{document}
